\section{Repository audit and the actual missing layer}
\label{sec:audit}

\subsection{Inspected baseline}

The audit was performed against the public \texttt{main} branch of
\texttt{VladimirReshetnikov/ProveIt}, whose head at the time of inspection was
commit \path{9df5cbdf39e6220be823f611987bd584d37711fd}.
The target directory's own index reports seventy-two delivered proof packets,
separated from the kernel-checked \texttt{ConciseIPRA} library.  The index is
commendably explicit that a packet may contain a complete human argument,
finite computations, source-level Lean, a compiled semantic lemma, an object
\(\Prf\;\PRAAx\) theorem, or only a target proposition; these statuses must not
be conflated.

The repository already contains much of the right mathematical vocabulary:
finite observations, precision contracts, backward demands, certificate
verification, bounded traces, coherent towers, ranked descent, relational
witnesses, stable observables, finite nets, height/list conversion, and
several application-specific algebraic and analytic estimates.  It also
contains a mature deeply embedded arithmetic backend.  The task is therefore
not to invent yet another coding of finite sequences or rational numbers.  It
is to organize these components around the one conclusion that matters for an
internal iPRA theorem.

\subsection{The backend that should be reused}

The \texttt{ConciseIPRA} library formalizes three mutually related systems:
the native scalar core, an intermediate equational theory, and ordinary iPRA.
For the present framework the ordinary interface is the most direct.  Its
principal declarations are:

\begin{itemize}
\item \texttt{PR k}: fixed primitive-recursive function symbols of arity
\(k\), generated by zero, successor, projection, composition, and primitive
recursion;
\item \texttt{PTm n}: terms with \(n\) scoped variables;
\item \texttt{Fm PTm n}: first-order formulas with equality, intuitionistic
connectives, and quantifiers;
\item \texttt{PRAAx}: defining equations, successor separation and
injectivity, and quantifier-free induction;
\item \texttt{Prf PRAAx Gamma phi}: the actual deeply embedded derivability
judgment;
\item structural proof constructors for weakening, substitution, cut,
quantifier rules, and equality transport;
\item \texttt{PRA.qfInd}: the object-theory quantifier-free induction rule.
\end{itemize}

This API is sufficient for a compiler.  It is also restrictive in exactly the
right way.  A Lean theorem that a natural-number function returns the expected
value is not automatically a \texttt{Prf} object.  A primitive-recursive
symbol's defining equations are available, but the arithmetic invariant still
has to be stated as a quantifier-free formula and proved with the embedded
rules.

The strongest reusable checked result found in the proof packets is the
local-to-global trace rule in
\path{foundations/ipra_local_global/lean/ObjectIteration.lean}.  In simplified
mathematical notation it says the following.

\begin{theorem}[Existing object trace rule, repository form]
\label{thm:existing-trace}
Let \(f\) be a \(k\)-ary primitive-recursive initializer, \(g\) a
\((k+2)\)-ary step function, and \(P(s,t,\vec x)\) a quantifier-free formula.
From embedded derivations of
\[
P(f(\vec x),0,\vec x)
\quad\text{and}\quad
P(s,t,\vec x)\longrightarrow P(g(t,s,\vec x),t+1,\vec x),
\]
with the exact scoped contexts prescribed by the library, the theorem
constructs an embedded derivation of
\[
\forall t\;P(\operatorname{run}_{f,g}(t,\vec x),t,\vec x).
\]
It also supplies a theorem specializing the universal conclusion at any
primitive-recursive fuel term.
\end{theorem}

This theorem is not a semantic induction argument disguised as a formalization
plan.  Its conclusion is the actual \texttt{Prf PRAAx} type.  The packet also
contains a closed counting example that exercises the recursion equations.
The stable-observables packet provides similarly concrete proof combinators:
a bounded invariant specialized at a fuel, extraction of a certificate from a
terminal invariant, comparison of two independently produced certificates,
and an end-to-end two-loop stability rule.

These are the correct foundations.  The PCOG proposal treats them as the
trusted reusable leaves and focuses on compiling a larger finite architecture
to those leaves.

\subsection{What the prior general frameworks achieved}

The foundation packets converge on a common core.

\begin{enumerate}
\item \emph{Certified finitary presentations} expose fixed codes,
verification predicates, restriction maps, and precision transformers.
\item \emph{Finite-observation compilers} replace unbounded mathematical
objects by finite tables and bounded runs.
\item \emph{Contract DAGs} propagate requested precision backward and execute
producers forward.
\item \emph{Relational finite use} keeps all valid answers when exact selection
would be unstable.
\item \emph{Coherent finite methods} prove compatibility under restriction and
use filtered lifting, ranked descent, and norm certificates.
\item \emph{Stable observables} compare compressed outputs of unrelated
certificates without requiring the witnesses themselves to agree.
\end{enumerate}

Several packets contain checked host-level finite lemmas, typed target
propositions, or a handful of embedded derivation combinators.  The unified
finitary-relational packet states the right generic theorem catalogue, but its
own ledger marks the generic \texttt{Prf PRAAx} compiler as an open engineering
task.  The contract-DAG packet likewise records that its generic trace and DAG
theorems are not implemented.  The finite-observation packet explicitly
retains the boundary between compiled finite tables and full deep
object-language application proofs.

This report therefore does not claim novelty for backward demands, finite
viability, or certificate comparison individually.  Its contribution is a
more exact \emph{normal form for the missing proof compiler}: every recursive
component becomes one checked trace package; every acyclic dependency becomes
one proof-producing rule; and the whole theorem is a staged topological fold.
That formulation is strong enough to cover the earlier frameworks and simple
enough to be implemented directly against the current \texttt{Prf} API.

\subsection{Recurring failure modes}

A review of theorem packets reveals six recurring gaps.

\paragraph{Semantic truth in place of derivability.}
A function is implemented in Lean or Python and shown correct in the host
logic, but no term of \texttt{Prf PRAAx [] phi} is constructed.  This is the
most common gap.

\paragraph{A goal structure mistaken for a proof.}
A Lean structure names all desired propositions, or a theorem accepts an
\texttt{InternalProof} field, but no instance fills it.  Type-checking the goal
is valuable, yet it does not establish the theorem.

\paragraph{Fixed-stage proofs without a uniform invariant.}
For every externally chosen \(n\), one can verify the \(n\)-th computation,
but the family is not connected by a single quantifier-free step theorem.  A
trace compiler needs the uniform step.

\paragraph{An exact selector inserted at a collision.}
Roots, eigenvectors, branches, signs, axes, or minimizers are chosen where only
a finite relation is justified.  This can encode LLPO or another nonconstructive
principle.  The repair is relational output plus a separate stability theorem.

\paragraph{A finite cardinality bound used as an enumeration bound.}
Knowing that a solution set has at most \(N\) elements gives no primitive-
recursive height cutoff.  A search compiler requires a cutoff, an explicit
cover, or a producer with a verifiable stopping certificate.

\paragraph{A universal interpreter smuggled into a uniform theorem.}
A framework quantifies over codes of arbitrary primitive-recursive programs
and assumes a total evaluator.  Diagonalization rules this out.  Fixed theorem
syntax must instead be compiled before entering arithmetic.

The PCOG discipline gives each failure an explicit type error: semantic facts
do not inhabit a rule's derivation field; absent uniformity cannot construct a
trace package; unstable choices remain clouds; cardinality is not accepted as
a bound; and graph syntax is staged rather than object-interpreted.

\subsection{The precise objective}

The desired compiler theorem should have the following shape.

\begin{quote}
Given a fixed finite graph; concrete primitive-recursive symbols for every
calculation; actual embedded proofs of every local deterministic,
relational, and stability contract; and actual embedded base and step proofs
for every recursive component, construct an embedded proof of the graph's
final exported formula.
\end{quote}

Notice what is absent: no appeal to the standard model, no proof-checker
reflection axiom, no arbitrary-program evaluator, no countable choice, and no
claim that a missing height bound can be obtained from a cardinality theorem.
The compiler is a theorem about composing \emph{proof objects already present
at the local interfaces}.  The rest of the report develops this statement in
detail.
